\documentclass[11pt, landscape]{article}
\usepackage[margin=0.75in]{geometry}
\usepackage[table, dvipsnames]{xcolor}
\usepackage{tabularx}
\usepackage{fontspec}
\usepackage{etoolbox}
\usepackage{enumitem}
\usepackage{environ}
\usepackage{xltabular}
\usepackage[most]{tcolorbox}

\setmainfont{PT Sans} 
\usepackage[portuguese]{babel}

% --- Color Palette ---
\definecolor{primaryBlue}{HTML}{2C3E50}
\definecolor{normalGreen}{HTML}{27AE60}
\definecolor{altOrange}{HTML}{E67E22}
\definecolor{excRed}{HTML}{C0392B}
\definecolor{lightGray}{HTML}{F8F9FA}
\definecolor{borderGray}{HTML}{D1D1D1}

% --- Table Styling ---
\renewcommand{\arraystretch}{1.6}
\setlength{\arrayrulewidth}{0.5pt}
\arrayrulecolor{borderGray}

% --- Counters ---
\newcounter{stepcount}
\newcounter{substepcount}

% --- Commands ---

\newcommand{\ucHeader}[5]{
	\noindent {\color{primaryBlue}\Huge \textbf{Caso de Uso: #1}}
	
	\vspace{1.5em}
	\noindent
	\begin{tabularx}{\textwidth}{|l|X|}
		\hline
		\rowcolor{lightGray} \textbf{Descrição} & #2 \\ \hline
		\textbf{Cenários} & #3 \\ \hline
		\rowcolor{lightGray} \textbf{Pré-condições} & #4 \\ \hline
		\textbf{Pós-condições} & #5 \\ \hline
	\end{tabularx}
}

\NewEnviron{flowtable}{
	\noindent
	\begin{xltabular}{\textwidth}{|>{\hsize=1.4\hsize}X|>{\hsize=0.8\hsize}X|>{\hsize=1.2\hsize}X|>{\hsize=0.6\hsize}X|}
		\hline
		\rowcolor{lightGray} 
		\textbf{Passos} & \textbf{Lógica de Negócio} & \textbf{Definição da API} & \textbf{Subsistema} \\ \hline
		\endfirsthead
		
		\hline
		\rowcolor{lightGray} 
		\textbf{Passos} & \textbf{Lógica de Negócio} & \textbf{Definição da API} & \textbf{Subsistema} \\ \hline
		\endhead
		
		\BODY
	\end{xltabular}
}

\newcommand{\step}[4]{
	\stepcounter{stepcount}
	\thestepcount. #1 & #2 & \texttt{#3} & #4 \\ \hline
}

\newcommand{\substep}[4]{
	\stepcounter{substepcount}
	\theparentstep.\thesubstepcount\ #1 & #2 & \texttt{#3} & #4 \\ \hline
}

% lazy step: use it when the operation will not be reflected in the system
\newcommand{\lstep}[1]{
	\step{#1}{-}{-}{-}
}

% lazy sub step: use it when the operation will not be reflected in the system
\newcommand{\lsubstep}[1]{
	\substep{#1}{-}{-}{-}
}

\newcommand{\normalFlowHeader}{
	\setcounter{stepcount}{0}
	
	\vspace{1em}
	
	\noindent
	\begin{tcolorbox}[
		enhanced,
		sharp corners,
		boxrule=0pt,
		frame hidden,
		colback=normalGreen!8!white,
		left=5pt, right=5pt,
		top=8pt, bottom=8pt,
		after skip=0pt,
		]
		{\color{normalGreen}\Large\bfseries Fluxo Normal}
	\end{tcolorbox}
	\nointerlineskip
}

\newcommand{\subFlowHeader}[4]{
	\setcounter{substepcount}{0}
	\def\theparentstep{#3}
	
	\noindent
	\begin{tcolorbox}[
		enhanced,
		sharp corners,
		boxrule=0pt,
		frame hidden,
		colback=#1!8!white,
		left=5pt, right=5pt,
		top=5pt, bottom=5pt,
		after skip=0pt,
		]
		{\color{#1}\Large\bfseries #2} \vspace{0.3em} \\
		\noindent\hspace{2pt}\vrule width 2pt depth 0.8ex height 2.2ex \hspace{0.5em}
		\begin{minipage}{\dimexpr\textwidth-25pt}
			\small \textbf{PASSO #3} \quad \textbullet \quad \textbf{CONDIÇÃO:} \textit{#4}
		\end{minipage}
	\end{tcolorbox}
	\nointerlineskip
}

\begin{document}

	\ucHeader{Iniciar Sessão}
	{Utilizador autentica-se na rede social}
	{Cenário 1}
	{Utilizador não autenticado}
	{Utilizador autenticado e com acesso à rede social}
	
	\normalFlowHeader
	\begin{flowtable}
		\lstep{Utilizador apresenta o seu identificador e senha}
		\step{Sistema valida identificador e senha}{Validar credenciais do utilizador}{validarUtilizador(id : String, senha : Strring) : boolean}{Utilizador}
		\step{Sistema apresenta uma lista de publicações feitas ou partilhadas pelas pessoas que segue}{Apresentar publicações de seguidores}{getPublicacoesSeguidores() : List<Publicacao>}{Publicacoes}
	\end{flowtable}
	
	\subFlowHeader{excRed}{Fluxo de Exceção}{2}{Credenciais inválidas}
	\begin{flowtable}
		\lstep{Sistema informa que credenciais são inválidas}
	\end{flowtable}

	\newpage
	
	\ucHeader{Adicionar Publicação}
	{Utilizador adiciona uma publicação}
	{Cenário 2}
	{Utilizador autenticado}
	{Sistema tem mais uma publicação}
	
	\normalFlowHeader
	\begin{flowtable}
		\lstep{Utilizador adiciona o conteúdo da publicação}
		\step{Sistema valida conteúdo da publicação}{Validar conteúdo de publicação}{validarConteudo(conteudo : String) : boolean}{Publicacoes}
		\step{Sistema adiciona a nova publicação}{Adicionar publicacao ao sistema}{adicionarPublicacao(conteudo : String) : void}{Publicacoes}
	\end{flowtable}
	
	\subFlowHeader{excRed}{Fluxo de Exceção}{2}{Conteúdo da publicação inválido}
	\begin{flowtable}
		\lsubstep{Sistema informa que conteúdo da publicação quebra as regras}
	\end{flowtable}

	\newpage
	
	\ucHeader{Comentar Publicação}
	{Utilizador comenta uma publicação de outro Utilizador}
	{Cenário 1}
	{Utilizador autenticado}
	{Comentário adicionado a publicação}
	
	\normalFlowHeader
	\begin{flowtable}
		\lstep{Utilizador seleciona uma publicação}
		\lstep{Utilizador escreve comentário}
		\step{Sistema valida conteúdo do comentário}{Validar conteúdo do comentário}{validarConteudo(conteudo : String) : boolean}{Publicacoes}
		\step{Sistema adiciona comentário à publicação}{Comentar publicação}{comentarPublicacao(comentario : String, codPubl : String) : void}{Publicacoes}
	\end{flowtable}
	
	\subFlowHeader{altOrange}{Fluxo Alternativo}{2}{Utilizador reage com emoji}
	\begin{flowtable}
		\lsubstep{Utilizador escolher emoji}
		\lsubstep{Regressa a 4}
	\end{flowtable}
	
	\newpage
	
	\subFlowHeader{excRed}{Fluxo de Exceção}{3}{Conteúdo do comentário inválido}
	\begin{flowtable}
		\lsubstep{Sistema informa que conteúdo do comentário quebras as regras}
	\end{flowtable}
	
\end{document}